% Capitulo: Do pipeline ao dashboard
\chapter{Do pipeline ao dashboard}
\label{cap:pipeline-dashboard}

Em muitas empresas, há uma parede entre quem constrói os dados e quem constrói os painéis. O
time de engenharia entrega tabelas e considera o trabalho terminado; o time de visualização
recebe as tabelas, aplica as próprias regras dentro da ferramenta de BI e publica. Quando um
número sai errado, cada lado aponta para o outro. Este capítulo trata de como derrubar essa
parede: definir as métricas uma vez, deixar claro quem é dono de quê e tratar o caminho do
dado bruto até o painel como um fluxo só.

\section{Métricas definidas uma vez}
\label{sec:camada-semantica}

O problema dos \enquote{mil dashboards e nenhuma verdade} (\veja{sec:mil-dashboards}) quase
sempre tem a mesma origem: a mesma métrica definida em vários lugares. O faturamento é
calculado num painel com um filtro, em outro com outro, e numa planilha com um terceiro
critério. Cada definição faz sentido para quem a escreveu, e juntas elas destroem a confiança
nos números.

A solução é uma camada semântica: um lugar único onde cada métrica é definida, com a fórmula,
os filtros padrão, o grão em que faz sentido e o dono, e de onde todos os painéis e
ferramentas leem. Quando a definição muda, muda num lugar só, e todos os consumidores passam a
usar a definição nova ao mesmo tempo. Como o capítulo~\ref{cap:nuvem} observou, a ideia não é
nova (\veja{sec:nomes-novos}); o que mudou foi que hoje ela pode viver junto com o código de
transformação, versionada e testada como ele.

A camada semântica ganhou um consumidor que torna a sua existência ainda mais importante: a
IA. Um assistente que responde perguntas sobre os dados precisa de definições claras para não
inventar a própria fórmula de faturamento (\veja{sec:metadados-contexto}). A mesma definição
que garante que dois painéis mostrem o mesmo número garante que o assistente responda com
ele.

\section{Quem é dono de quê}
\label{sec:donos}

A parede entre engenharia e visualização é, em boa parte, um problema de responsabilidade mal
definida. Quando ninguém sabe quem é dono da métrica, todos acham que é o outro. A
Tabela~\ref{tab:donos} propõe uma divisão que costuma funcionar.

\begin{table}[htbp]
  \centering
  \caption{Uma divisão de responsabilidades do dado de origem ao painel.}
  \label{tab:donos}
  \small
  \begin{tabularx}{\textwidth}{@{}>{\raggedright\arraybackslash}p{0.2\textwidth}
      >{\raggedright\arraybackslash}p{0.28\textwidth}>{\raggedright\arraybackslash}X@{}}
    \toprule
    O quê & Dono & Responde por \\
    \midrule
    Dado na origem & Área dona do sistema de origem & Registrar corretamente, avisar mudanças no sistema \\
    \addlinespace
    Plataforma & Engenharia de dados & Ingestão, disponibilidade, desempenho, custo \\
    \addlinespace
    Modelos de dados & Analytics engineering & Transformação correta, testes, atualização no prazo \\
    \addlinespace
    Métrica & Dono de negócio da métrica & A definição: o que entra, o que sai, em que critério \\
    \addlinespace
    Painel & Quem o construiu, com a área usuária & Desenho, uso, manutenção e desligamento \\
    \bottomrule
  \end{tabularx}
  \fonte{elaborado pelo autor.}
\end{table}

Duas linhas da tabela costumam surpreender. A primeira é que o dono da métrica é alguém do
negócio, e não do time de dados. O time de dados implementa a definição e garante que ela seja
calculada corretamente, mas quem decide se a devolução desconta do faturamento é a área que
usa o número para decidir. A segunda é que todo painel tem um dono, inclusive para decidir
quando ele deixa de existir. Painéis sem dono são os que se acumulam, consomem atualização e
ninguém tem coragem de apagar (capítulo~\ref{cap:frequencia}).

\section{Um fluxo só}
\label{sec:fluxo-unico}

Com as métricas centralizadas e os donos definidos, o passo seguinte é tratar o caminho do dado
até o painel como um fluxo único, e não como duas entregas separadas.

Na prática, isso significa algumas coisas. A primeira é manter a lógica fora do painel: regras
de negócio, filtros padrão e cálculos ficam na camada de transformação ou na camada semântica,
onde são versionados e testados; o painel só apresenta. A segunda é conhecer a linhagem até o
fim: quando um modelo vai ser alterado, o time sabe quais painéis dependem dele e pode avaliar
o impacto antes, em vez de descobrir depois, por uma reclamação. A terceira, quando a
ferramenta permite, é versionar e publicar os painéis com o mesmo processo dos modelos
(capítulo~\ref{cap:ci-cd}). E a quarta é organizacional: um quadro de trabalho único para
engenharia e visualização, em que uma demanda de painel e a demanda de dados que ela depende
aparecem juntas, como partes da mesma entrega (Parte~\ref{parte:agil}).

O efeito mais visível dessa integração é na conversa com o negócio. Em vez de o pedido passar
por dois times, com duas filas e duas interpretações, ele é tratado do começo ao fim por um
fluxo só, com um responsável pela entrega. O efeito menos visível, e talvez mais importante, é
na qualidade: quando quem constrói o painel enxerga o modelo, e quem constrói o modelo
enxerga o painel, os erros aparecem antes de chegar à reunião.

\rascunho{
  Relato: a integração entre o time de engenharia e o de visualização num projeto real (como
  era antes, o que mudou na forma de trabalhar, o que melhorou e o que ainda é difícil).
}

\begin{leituras}
  \leitura{kimball2013}{Para a base de tudo isto: métricas e dimensões conformadas, definidas
    uma vez e compartilhadas por todos os processos.}
  \leitura{skelton2019}{Para pensar a divisão de responsabilidades entre times de plataforma
    e times que entregam valor ao negócio.}
\end{leituras}
